\documentclass[11pt]{article}
\usepackage[T1]{fontenc}
\usepackage{lmodern,amsmath,amssymb,amsthm,mathtools,microtype}
\usepackage[margin=1in]{geometry}
\usepackage[colorlinks=true,allcolors=blue]{hyperref}
\hypersetup{pdftitle={A negative-curvature normal barrier on an algebraic cone that is not a hyperbolicity cone},pdfauthor={Alec Kriebel}}
\newtheorem{theorem}{Theorem}
\newcommand{\R}{\mathbb R}
\newcommand{\inter}{\operatorname{int}}
\newcommand{\diag}{\operatorname{diag}}
\title{A negative-curvature normal barrier on an algebraic cone\protect\\that is not a hyperbolicity cone}
\author{Alec Kriebel\thanks{Independent researcher. ORCID: \href{https://orcid.org/0009-0001-9320-500X}{0009-0001-9320-500X}.}}
\date{October 17, 2026 \quad Version 1.0.0}
\begin{document}
\maketitle
\begin{abstract}
We give an explicit affirmative answer to Tun\c{c}el's 2017 existential
question about algebraic convex cones with negative-curvature normal
barriers. The one-sided power cone
$K=\{(x,y,z)\geq0:xy^2\geq z^3\}$ admits the barrier
$F=-25\log\bigl(z(xy^2-z^3)\bigr)$ with parameter $100$ and
$D^3F(a)[h,u,u]\leq0$ for every $a\in\inter K$, $h\in K$, and
$u\in\R^3$. Exact polynomial and positive-semidefinite certificates prove
the curvature inequality, including its degenerate directions. A boundary
divisibility argument rules out every homogeneous hyperbolic polynomial
defining $K$ in $\R^3$. The cone itself is a known member of Hildebrand's
classification; the contribution here is the explicit barrier certificate
and its verification for the stated existential problem.
\end{abstract}

\section{The problem and the example}
Tun\c{c}el \cite[p.~799, Open Problem~1]{T} asks whether an algebraic convex
cone can admit a normal barrier with negative curvature without being a
hyperbolicity cone. Here algebraic means defined by homogeneous polynomial
inequalities. For a regular cone (closed, pointed, convex, with nonempty
interior), negative curvature means
\begin{equation}\label{negative}
 D^3F(a)[h,u,u]\leq0 \quad(a\in\inter K,\ h\in K,\ u\in\R^3).
\end{equation}
This is the third-derivative convention of \cite{T,NT}, equivalently the
concavity of $a\mapsto DF(a)[h]$ for fixed $h\in K$ \cite[Definition~3.1]{DTV}.
A normal barrier is a logarithmically homogeneous self-concordant barrier;
we also verify its gradient bound directly.

\begin{theorem}\label{main}
Let
\[
 K=\{(x,y,z)\in\R^3:x,y,z\geq0,\ p(x,y,z):=xy^2-z^3\geq0\}.
\]
Then $K$ is a regular algebraic convex cone. On $\inter K$, the function
\[
 f=-\log p-\log z,\qquad F=25f
\]
is a normal barrier with parameter $100$ satisfying \eqref{negative}.
No homogeneous polynomial in three variables has $K$ as its
hyperbolicity cone.
\end{theorem}

The underlying cone appears in Hildebrand \cite[Theorem~3.2, family~5]{H}:
take the conjugate exponents $3,3/2$ and coordinates
$(x_1,x_2,x_3)=(x,z,y)$. Self-concordant barriers for this same nonnegative
power-cone family are also known: Roy--Xiao \cite[Theorem~2]{RX} give,
for this cone, the parameter-$3$ barrier
\[
 B=-\log(x^{1/3}y^{2/3}-z)-\tfrac23\log x-\tfrac13\log y-\log z.
\]
Our parameter $100$ is not an improvement of their self-concordance
parameter. The property established here is the negative-curvature
certificate for the displayed polynomial-logarithmic barrier, together
with the all-polynomial hyperbolicity obstruction. The theorem concerns this
one-sided cone and the original existential question. It does not resolve
the broader question whether every regular cone admits such a barrier
\cite[Section~6]{DTV}. All arguments below are analytic; the accompanying
exact computations are independently checkable supplements.

\section{Geometry and normality}
The displayed homogeneous inequalities make $K$ closed and conic. The
nonnegative orthant makes it pointed; $(2,1,1)$ is strictly feasible, and
\[
 \inter K=\{x,y,z>0:xy^2>z^3\}.
\]
For $m(x,y)=x^{1/3}y^{2/3}$ on the positive orthant,
\[
 \nabla^2m=-\frac29x^{-5/3}y^{-4/3}
 \begin{pmatrix}y^2&-xy\\-xy&x^2\end{pmatrix}\preceq0.
\]
Concavity extends to the closed orthant by continuity. Thus
$K=\{x,y\geq0,\ 0\leq z\leq m(x,y)\}$ is convex.

Put $g=z^3/y^2$ and $\delta=x-g>0$. Then
$f=-\log\delta-2\log y-\log z$. At a fixed base point, for any direction
$u$, define the linear functionals
\[
 b_u=u_y/y,\quad c_u=u_z/z,\quad
 w_u=\frac{u_x-g(3c_u-2b_u)}{\delta},\quad
 v_u=\sqrt{6g/\delta}\,(c_u-b_u),\quad \ell_u=c_u-4b_u.
\]
Direct differentiation gives
\begin{align}
 \|u\|_f^2:=D^2f[u,u]&=w_u^2+v_u^2+2b_u^2+c_u^2,\label{hess}\\
 D^3f[u,u,u]&=-2w_u^3-3w_uv_u^2+\ell_uv_u^2-4b_u^3-2c_u^3.\label{cubic}
\end{align}
For example, $Dg[u]=g(3c_u-2b_u)$,
$D^2g[u,u]=6g(c_u-b_u)^2$, and
$D^3g[u,u,u]=6g(c_u-b_u)^2(c_u-4b_u)$, which yield these identities.
Equation \eqref{hess} is positive definite: zero forces $b_u=c_u=w_u=0$
and hence $u=0$.

For completeness we prove the mixed form of self-concordance used in
\cite[Section~2.2]{DTV}. Polarizing \eqref{cubic} gives
\begin{align*}
 D^3f[h,u,u]={}&-2w_hw_u^2-w_hv_u^2-2w_uv_hv_u\\
 &+\tfrac13(\ell_hv_u^2+2\ell_uv_hv_u)-4b_hb_u^2-2c_hc_u^2.
\end{align*}
Weighted Cauchy--Schwarz gives $|\ell_a|\leq3\sqrt{2b_a^2+c_a^2}
\leq3\|a\|_f$. The four groups in this display have absolute values
bounded respectively by $2,3,3,2$ times $\|h\|_f\|u\|_f^2$.
For the last bound, use
\[
 |4b_hb_u^2+2c_hc_u^2|
 \leq2\sqrt{2b_h^2+c_h^2}\sqrt{2b_u^4+c_u^4}
 \leq2\|h\|_f\|u\|_f^2.
\]
Consequently $|D^3f[h,u,u]|\leq10\|h\|_f\|u\|_f^2$. Since
$\|u\|_F=5\|u\|_f$, scaling by $25$ gives
\[
 |D^3F[h,u,u]|\leq2\|h\|_F\|u\|_F^2,
\]
including the standard scalar self-concordance inequality when $h=u$.

The polynomial $zp$ has degree four, so
$F(ta)=F(a)-100\log t$. Differentiation yields
$\nabla^2F(a)a=-\nabla F(a)$ and $a^T\nabla^2F(a)a=100$; hence
\[
 \nabla F(a)^T\nabla^2F(a)^{-1}\nabla F(a)=100.
\]
At every finite boundary point of $K$, either $z=0$ or $p=0$.
Thus along every interior sequence tending to that point, $zp\to0^+$
and $F\to+\infty$. Extended by $+\infty$ outside $\inter K$, $F$ is a
closed convex function. These facts prove its claimed normality.

\section{The negative-curvature certificate}
The diagonal maps $D=\diag(\alpha,\beta,\gamma)$ with positive entries
and $\alpha\beta^2=\gamma^3$ are automorphisms of $K$, and
$f(Dq)=f(q)-4\log\gamma$. For a fixed interior point $(x,y,z)$, take
$D=\diag(y^2/z^3,1/y,1/z)$. It sends that point to $(r,1,1)$,
where $r=xy^2/z^3>1$. Differentiating with this $D$ held fixed gives
\[
 D^3f(q)[h,u,u]=D^3f(Dq)[Dh,Du,Du].
\]
It is therefore enough to prove the matrix inequality at $(r,1,1)$
for every $r>1$.

Every $h\in K$ with $h_y>0$ has the decomposition
\[
 h=h_y(t^3,1,t)+\bigl(h_x-h_z^3/h_y^2\bigr)e_x,\qquad t=h_z/h_y\geq0,
\]
with both coefficients nonnegative. If $h_y=0$, then $h_z=0$ and $h$
is a nonnegative multiple of $e_x$. By linearity in $h$, it remains to
check $h(t)=(t^3,1,t)$ for $t\geq0$, and $e_x$.

Set $s=r-1>0$ and $M=-s^3D^3f(r,1,1)[h(t)]$. Its symmetric entries are
\begin{align*}
 M_{11}&=2(t-1)^2(t+2),\\
 M_{12}&=2(t-1)(2t^2+2t-1-3r),\\
 M_{13}&=6(t-1)(r-t^2-t+1),\\
 M_{22}&=4r^3+12r^2-18r^2t+6rt^3-6rt+2t^3,\\
 M_{23}&=12r^2t-18r^2-6rt^3+24rt-6r-6t^3,\\
 M_{33}&=6(r+2)t^3+2(r^3-6r^2-18r-4)t+12r(r+2).
\end{align*}
These follow, for instance, by differentiating
$f_{ij}=p_ip_j/p^2-p_{ij}/p+\mathbf1_{i=j=z}/z^2$.
The next two leading minors are the following polynomial identities:
\begin{align}
 M_{11}M_{22}-M_{12}^2&=4s(t-1)^2P,\label{minor}\\
 \det M&=8s^4(t-1)^2R,\label{det}
\end{align}
where
\begin{align*}
 P&=2(t+2)s^2-3(t-1)(3t+5)s+3(t-1)^2(t+2)^2,\\
 R&=(t+2)\bigl(2ts^2-3(t-1)(3t+1)s
             +3(t-1)^2(t^2+4t+1)\bigr).
\end{align*}
For $0\leq t<1$, every term in these expressions is nonnegative and
the constant terms are positive. At $t=1$, $P=R=6s^2>0$.
For $t\geq1$ use the exact square identities
\begin{align}
 8(t+2)P={}&[4(t+2)s-3(t-1)(3t+5)]^2\notag\\
          &+3(t-1)^2(8t^3+21t^2+6t-11),\label{squarep}\\
 \frac{8tR}{t+2}={}&[4ts-3(t-1)(3t+1)]^2\notag\\
          &+3(t-1)^3(8t^2+13t+3).\label{squarer}
\end{align}
The first remainder polynomial, at $t=1+v$, equals
$8v^3+45v^2+72v+24>0$ for $v\geq0$. Both remainders are strictly
positive for $t>1$, and all divisors used are positive. Hence $P,R>0$
throughout $s>0,t\geq0$. If $t\ne1$, the first three leading minors
are positive, so Sylvester's criterion gives $M\succ0$.

At the exceptional direction $t=1$, substitution instead gives
\[
 M=2s^2\begin{pmatrix}0&0&0\\0&2r+1&-3\\0&-3&r+2\end{pmatrix}\succeq0,
\]
because the lower block has positive diagonal and determinant
$(2r+1)(r+2)-9=(r-1)(2r+7)>0$.
For the remaining endpoint there is a direct Gram certificate:
\[
 -s^3D^3f(r,1,1)[e_x]=2aa^T+6sbb^T\succeq0,\qquad
 a=(1,2,-3)^T,\quad b=(0,1,-1)^T.
\]
Thus all cone directions are covered. Congruence under $D$, linearity
in $h$, and multiplication by $25$ prove \eqref{negative}.

\section{Excluding every hyperbolicity polynomial}
Recall that a homogeneous polynomial $H$ is hyperbolic in $e$ if
$H(e)>0$ and $H(v+te)$ has only real roots for every $v$.
Its closed hyperbolicity cone is the closure of the positive-eigenvalue
region in the convention of \cite{T}. The direction $e$ lies in the
interior of that cone. Boundary eigenvalues are nonnegative by continuity
from the positive region. The boundary lies in $\{H=0\}$: if all eigenvalues
there were positive, continuity of roots (with fixed
nonzero leading coefficient) would give an interior neighborhood.

Suppose that this cone were $K$. Then $e\in\inter K$, and the curved
boundary patch implies
\[
 H(z^3/y^2,y,z)=0\qquad(y,z>0).
\]
Clearing denominators gives a polynomial identity on an open set in
$\R^2$, hence an identity everywhere. In $\R(y,z)[x]$ this says
$x-z^3/y^2$ divides $H$. The polynomial $p=y^2x-z^3$ is primitive
over $\R[y,z]$, since $\gcd(y^2,z^3)=1$, and is irreducible over the
fraction field, being linear in $x$. Gauss's lemma gives $p\mid H$
in $\R[x,y,z]$.

Write $a=(e_xe_y^2)^{1/3}$ and $b=e_z$, so $a>b>0$. For $v=(0,0,1)$,
\begin{align*}
 p(v+te)&=a^3t^3-(1+bt)^3\\
 &=((a-b)t-1)\bigl((a^2+ab+b^2)t^2+(a+2b)t+1\bigr).
\end{align*}
The quadratic discriminant is $(a+2b)^2-4(a^2+ab+b^2)=-3a^2<0$.
Therefore $p(v+te)$ has nonreal roots. Since $p\mid H$ and $H(e)>0$,
the restricted product is nonzero and these roots cannot cancel.
This contradicts hyperbolicity in $e$, proving the final assertion of
Theorem~\ref{main}. The argument treats arbitrary homogeneous $H$,
without an assumed degree or a particular defining polynomial.

\section*{Verification, attribution, and research status}
The supplement supplies independently authored exact rational-polynomial
checkers for the derivative, minor, square, and endpoint identities, with
additional logarithm-series checks and deliberately false controls.
The proof above establishes the universal inequalities; finite samples
are supplementary evidence. The term $-\log z$ is substantive: for
$f_0=-\log p$, the matrix $-s^3D^3f_0[h]$ has determinant $-24576$
at base $(5,1,1)$ and direction $h=(8,1,2)$.

The cone's prior occurrence in \cite{H} is explicitly credited; its
hyperbolic affine spheres are distinct from polynomial hyperbolicity.
Definitions and related barrier results are credited to \cite{NT,DTV}.
A bounded primary-source priority audit accompanies the supplement;
no assertion of exhaustive literature coverage or first-ever priority
is made.

\textbf{AI and review disclosure.} AI tools were used extensively in
developing, checking, and writing this result and its verification
materials. This preprint is unrefereed. The accompanying independent
adversarial AI reviews do not constitute external human peer review
or proof-assistant certification.

\begin{thebibliography}{9}
\bibitem{T} L.~Tun\c{c}el, \emph{Polynomial optimization with a focus on
hyperbolic polynomials}, Oberwolfach Report 14/2017, 797--800,
Open Problem~1 on p.~799.
\href{https://doi.org/10.4171/owr/2017/14}{doi:10.4171/owr/2017/14}.
\bibitem{NT} Yu.~Nesterov and L.~Tun\c{c}el,
\emph{Local superlinear convergence of polynomial-time interior-point
methods for hyperbolic cone optimization problems}, SIAM J. Optim.
\textbf{26} (2016), 139--170; \href{https://arxiv.org/abs/1412.1857}{arXiv:1412.1857}.
\bibitem{DTV} J.~Dahl, L.~Tun\c{c}el, and L.~Vandenberghe,
\emph{New complexity bounds for primal--dual interior-point algorithms
in conic optimization}, \href{https://arxiv.org/abs/2509.10263}{arXiv:2509.10263v2}
(version record June 12, 2026). Sections 2.2, 3.1, and 6.
The inspected PDF is internally revised June 16; the HTML representation
instead says August 24. Both representations were checked.
\bibitem{H} R.~Hildebrand, \emph{Analytic formulas for complete hyperbolic
affine spheres}, \href{https://arxiv.org/abs/1305.4814}{arXiv:1305.4814},
Theorem~3.2, family~5.
\bibitem{RX} S.~Roy and L.~Xiao, \emph{On self-concordant barriers for
generalized power cones}, Microsoft Research Technical Report
MSR-TR-2018-3, January 30, 2018, Theorem~2.
\href{https://optimization-online.org/2018/01/6441/}{Optimization Online}.
\end{thebibliography}
\end{document}
